\documentclass[oneside,english]{amsart}
\usepackage[T1]{fontenc}
\usepackage[latin9]{inputenc}
\usepackage{babel}
\usepackage{mathtools}
\usepackage{amstext}
\usepackage{amsthm}
\usepackage{amssymb}
\usepackage{xr-hyper}
\usepackage[pdfusetitle,
 bookmarks=true,bookmarksnumbered=false,bookmarksopen=false,
 breaklinks=false,pdfborder={0 0 1},backref=false,colorlinks=false]
 {hyperref}
% Cross-references to Papers I and II: build ../theory and ../tool first.
\externaldocument[I-]{../theory/crosscheck-theory}
\externaldocument[II-]{../tool/crosscheck-tool}

\makeatletter
\numberwithin{equation}{section}
\theoremstyle{plain}
\newtheorem{thm}{\protect\theoremname}
\theoremstyle{definition}
\newtheorem{defn}[thm]{\protect\definitionname}
\theoremstyle{plain}
\newtheorem{prop}[thm]{\protect\propositionname}
\theoremstyle{plain}
\newtheorem{lem}[thm]{\protect\lemmaname}
\theoremstyle{remark}
\newtheorem*{rem*}{\protect\remarkname}

\newcommand{\lean}[1]{{\ttfamily\small #1}}
\emergencystretch=3em
\makeatother

\usepackage{booktabs}
\usepackage{array}
\providecommand{\definitionname}{Definition}
\providecommand{\lemmaname}{Lemma}
\providecommand{\propositionname}{Proposition}
\providecommand{\remarkname}{Remark}
\providecommand{\theoremname}{Theorem}

\begin{document}
\title[CrossCheck III: Certified Cycle Bounds of Machine Code]{CrossCheck III: Certified Cycle Bounds of Machine Code, Loop Bounds Included}
\author{Luca Dall'Olio}
\thanks{Alter Ego Engineering, 5 rue du Docteur Lequeux, 92330 Sceaux, France. Email: \texttt{ld@aego.ai}. ORCID: \href{https://orcid.org/0009-0006-8839-5909}{0009-0006-8839-5909}.}
\begin{abstract}
We certify worst-case cycle bounds of compiled machine code with a checker whose soundness is proved in Lean~4. A bound is certified by a dual solution of the implicit path enumeration program, potentials and multipliers checked with integer arithmetic on a control-flow graph; the checker builds the graph itself from the machine code, with a decoder and a cycle table of the Cortex-M0 written in Lean, so that the specification of the processor replaces the analysis of the binary in the trusted base. The loop bounds, which the implicit path enumeration needs and which come from the source program, can themselves be certified on the machine code: the certificate gives every instruction an invariant on the registers and on the operands of the operation that set the flags, and every loop a potential that decreases on its back edges, all checked by Farkas multipliers against a semantics of the instructions' effects on registers and flags. On the 29 TACLeBench kernels compiled with Clang for a Cortex-M0, all 116 functions that receive a cycle bound are certified, and all 116 loops of the run-time division routines and 143 of the 403 loops of the kernels' own code are bounded on the machine code with nothing trusted about them (200 of the kernels' loops when the checker also verifies, on the code of every function, the calling convention that makes callees preserve $r_4, \dots, r_7$ and the stack pointer, which it establishes in all 29 kernels, assuming only that no store through a pointer overwrites the words where registers are saved; 260 if, in addition, functions that take no address of their frame other than the frame pointer are assumed not to have it written by other code; 261 when the kernels are compiled without frame pointer); with both assumptions, 17 loop rows remain trusted in the 121 certified cycle bounds. No bound is exceeded in emulated runs. The other loops keep their source bounds, which remain trusted. Linked with LLVM's compiler-rt, whose floating-point routines are then bounded on their machine code too, 187 of the 267 functions of the kernels receive a certified cycle bound, the entry functions of 10 kernels among them. Jumps through a table, the code of a \texttt{switch}, are followed by evaluating their window on the code, which the checker repeats and Lean proves sound; on the 23 of TACLeBench's larger programs whose analysis finishes within an hour, 316 of 455 functions are certified. We also describe the cycle and stack bounds of the same binaries, their comparison with the WCET analyzer OTAWA, and a check of time budgets after compilation.
\end{abstract}

\maketitle

\section{Introduction}
\label{sec:intro}

A bound on the execution time of a real-time function is a property of its machine code on a processor. The companion papers derive bounds of source programs: Paper~I \cite{crosscheckTheory} gives the theory and a checker of cost certificates proved sound in Lean, and Paper~II \cite{crosscheckTool} a Clang plugin that bounds the loops of C and C++ functions and emits such certificates. This paper takes the step to the binary. The bounds of the source loops, matched with the loops of the compiled code through its debug information, give a cycle bound of each function of the binary on a processor whose instruction timings are fixed, by the timing schema or by implicit path enumeration (IPET) \cite{li1995ilp}. The question we address is what has to be trusted for such a bound, and how much of it a small verified checker can take over.

The chain from a source loop bound to a cycle bound of the binary has four links: the loop bound itself, certified on the source for the functions of Paper~II's fragment; the correspondence between source and binary loops, which assumes that the compiler preserves iteration counts; the control-flow graph of the binary and the computation of the bound; and the timing of the processor. We make the checker responsible for the third link (Sections~\ref{sec:graphs} and~\ref{sec:machine}): the certificate holds the machine code, the checker decodes it and builds the graph with a specification of the processor written in Lean, and checks a dual solution of the IPET program on it. For the second link we give certificates of loop bounds on the machine code itself (Section~\ref{sec:loops}): with them, a loop's bound no longer depends on the compiler, on the debug information nor on the source analysis. What remains trusted is the specification of the processor: the decoder, the cycle table, the effects of the instructions on registers and flags, and the absence of interrupts.

\subsection*{Contributions}
\begin{itemize}
  \item Cycle certificates for control-flow graphs with calls and loop rows, as dual solutions of the IPET program, and their soundness (Proposition~\ref{prop:cycles}; known in substance from Maroneze et al.\ \cite{maroneze2014formally}, here with calls and recursion by depth of runs).
  \item The check on machine code: a decoder and cycle table of the Cortex-M0 in Lean, the graph built by the checker, and the theorem that every run of the code is a walk of that graph (Theorem~\ref{thm:machine-cycles}).
  \item Loop bounds certified on the machine code: a data semantics of the instructions' effects on the low registers and the flags (Definition~\ref{def:data-run}), certificates of invariants and potentials (Definition~\ref{def:loop-certificate}), and the theorem that the rows of the certified loops hold on every run, so that they need not be trusted (Theorem~\ref{thm:machine-loops}); and the calling convention of the callees checked on their code instead of assumed, by an induction on the depth of calls in a semantics that links a callee's run to its caller's states (Theorem~\ref{thm:machine-conv}).
  \item An evaluation on the TACLeBench kernels compiled for a Cortex-M0 (Section~\ref{sec:evaluation}): cycle and stack bounds compared with emulated runs, certificates, loop bounds certified on the machine code, a comparison with OTAWA, and the check of time budgets after compilation.
\end{itemize}
All the statements of Sections~\ref{sec:graphs} to~\ref{sec:loops} are mechanized in Lean~4 without axioms beyond the three standard ones (Section~\ref{sec:mechanization}).

\section{Related Work}
\label{sec:related}

\paragraph{Implicit path enumeration.} IPET bounds the cost of a program by a linear program over the execution counts of the edges of its control-flow graph, with flow conservation and loop bounds as constraints \cite{li1995ilp}; it is the path analysis of industrial WCET analyzers \cite{wilhelm2008worst} and of OTAWA \cite{ballabriga2010otawa}. Constraints on execution counts beyond loop bounds go back to the markers, scopes and loop sequences of Puschner and Koza \cite{puschner1989calculating}. A feasible solution of the dual of the linear program bounds its optimum (weak duality \cite{schrijver1986theory}), which is what our certificates check.

\paragraph{Verified WCET estimation.} Blazy, Maroneze and Pichardie verified a loop-bound estimation in Coq on CompCert's RTL \cite{blazy2013formal}, and Maroneze et al.\ built on it a WCET estimation for CompCert, whose IPET program is generated from the semantics of CompCert's assembly language and whose solution is validated by a Farkas certificate of the infeasibility of a larger cost \cite{maroneze2014formally}. Their guarantees hold for code produced by CompCert; ours hold for any binary, relative to a specification of the processor, and the loop bounds are checked on the machine code instead of being proved for the compiler's intermediate language.

\paragraph{Loop bounds on binaries.} The aiT analyzer derives loop bounds on executables by a data-flow analysis of how loop counters change per iteration, start values from its value analysis, and the evaluation of the loop tests, with overflows of the counters taken into account; through the value analysis, the counters may be registers or memory cells, also modified through pointers, and the result is an interval of iteration counts, used after a faster pattern-based analysis \cite{cullmann2007data}. The bounds are computed, not certified. Congruences of integer variables were introduced as a static analysis by Granger, who also combines them with intervals by rounding each bound to the class \cite{granger1989static}; we use them, modulo powers of two, so that counters stepped by a constant that leave their loop on equality get a bound, and we check them in the certificate. For the seL4 kernel, Sewell, Kam and Heiser determine loop bounds of the binary with the translation-validation toolchain of the kernel's verification: SMT models of the binary's loops, by unrolling or by induction over a symbolic iteration, use facts proved of the C code, and a binary search refines the bound \cite{sewell2016complete,sewell2017high}. Their example of a binary loop is a register stepped by 4 that leaves on inequality with 256, bounded with the help of the divisibility of its initial value by 4, the case for which we use congruences; with the information of the binary alone, they bound 47 of the 67 loops of seL4, the others failing for lack of an aliasing analysis of memory and of invariants maintained by the loop's environment \cite{sewell2016complete}, which are also the main causes of our remaining failures (Section~\ref{sec:evaluation}). The proofs are those of SMT solvers inside their toolchain, and the control-flow graph is obtained by a decompiler that proves its output against the Cambridge model of the ARM instruction set \cite{fox2010trustworthy}, except for a final simplification, whose formalization they leave as future work. Our loop bounds are certificates of invariants and potentials checked by a small program proved sound in Lean; they are weaker (no memory, no C-level facts) and the specification of the processor is ours, not validated against hardware.

\paragraph{Certified resource bounds.} Proof-carrying code ships with a program a proof that a small trusted checker validates \cite{necula1997proof}. Crary and Weirich certified running-time bounds of typed assembly code with a decidable type system in which function types carry a virtual clock (a device they attribute to Necula and Lee), produced by a compiler from source programs annotated with running times \cite{crary2000resource}; time is counted in function calls in their calculus and in backward jumps in their typed assembly language, the checker is the type checker, and that well-typed programs keep within their bounds follows from type safety, proved on paper. Carbonneaux et al.\ derive resource bounds of C programs certified in Coq against CompCert's Clight \cite{carbonneaux2015compositional}. Our certificates are in processor cycles of the binary of an unmodified compiler, and the checker's soundness is a Lean theorem about the machine code.

\section{Cycle Certificates on Control-Flow Graphs}
\label{sec:graphs}

The loop bounds of the source program become a cycle bound of the compiled code through its control-flow graph: the implicit path enumeration technique (IPET) of Li and Malik \cite{li1995ilp} maximizes the cost of the edge counts that satisfy flow conservation and the loop bounds, a linear program. A feasible solution of its dual is a certificate of the bound, checked without the solver; this is weak duality of linear programming \cite{schrijver1986theory}. Maroneze et al.\ validated in Coq the solution of the IPET program of their WCET estimation with a Farkas certificate of the infeasibility of a larger cost \cite{maroneze2014formally}; the certificate below is the same idea, in the form of potentials and multipliers, with costs on edges and calls bounded by the claims of the callees.

\begin{defn}[Cycle certificate]
\label{def:cycle-certificate}
A function is a finite directed multigraph with an entry node $s$ and an exit node $t$; each edge $e = (u, v)$ has a weight $w_e \in \mathbb{N}$ and a list $\mathrm{calls}(e)$ of functions, and the function has finitely many \emph{rows} $\sum_e a_{ke}\, x_e \le b_k$ ($a_{ke}, b_k \in \mathbb{Z}$). A \emph{run} of depth $n + 1$ of a function is a walk $e_1 \cdots e_m$ from $s$ to $t$ together with, for each step $i$, runs of depth $n$ of the functions of $\mathrm{calls}(e_i)$ (there is no run of depth 0) such that the edge counts $x_e = |\{i : e_i = e\}|$ satisfy the rows; its cost is $\sum_i (w_{e_i} + c_i)$, where $c_i$ is the sum of the costs of the runs of step $i$. Given claimed bounds $B$ for all functions of a program, a \emph{certificate} of a function is an integer $D > 0$, integers $\pi_v$ (potentials) and natural numbers $\lambda_k$ (multipliers) such that
\[
  D\Bigl(w_e + \sum_{g \in \mathrm{calls}(e)} B(g)\Bigr) + \pi_v - \pi_u \;\le\; \sum_k \lambda_k\, a_{ke} \quad \text{for every edge } e = (u, v),
\]
and
\[
  \pi_s - \pi_t + \sum_k \lambda_k\, b_k \;\le\; D\, B(f).
\]
\end{defn}

\begin{prop}[Cycle bounds]
\label{prop:cycles}
If every function of a program has a certificate for the claimed bounds $B$, every run of every function $f$ costs at most $B(f)$.
\end{prop}

\begin{proof}
By induction on the depth of the run; the runs of the calls of step $i$ cost at most $\sum_{g \in \mathrm{calls}(e_i)} B(g)$ by the induction hypothesis. Summing the edge condition along the walk, the potentials telescope:
\[
  D \cdot \mathrm{cost} + \pi_t - \pi_s \;\le\; \sum_i \sum_k \lambda_k\, a_{k e_i} \;=\; \sum_k \lambda_k \sum_e a_{ke}\, x_e \;\le\; \sum_k \lambda_k\, b_k,
\]
the last step because the rows hold on the run and $\lambda_k \ge 0$. With the second condition, $D \cdot \mathrm{cost} \le D\, B(f)$.
\end{proof}

The induction is on the depth of a finite run, so the claimed bounds may be circular: the certificates of recursive functions are checked against the claims of each other, and the conclusion holds for every run that completes. The bounds of loops become rows: a loop entered by edges $E$, whose back edges $R$ may be taken at most $N$ times per entry, gives $\sum_{e \in R} x_e - N \sum_{e \in E} x_e \le 0$ (with right-hand side $N$ if the loop header is the entry), the sum of its per-entry bounds. The best certificate is an optimal solution of the dual of the IPET program; the bound it gives is the IPET bound, at most that of the timing schema \cite{puschner1989calculating} on the same graph. The executable check is the function \lean{Cycles.\allowbreak checkProgram} of the certificate checker.


\section{Machine Code}
\label{sec:machine}

The graph need not be given: it can be built from the machine code, so that the remaining interface to the binary is a specification of the processor. Sewell, Kam and Heiser obtained the control-flow graph of the seL4 binary from a decompiler that proves its output against a model of the ARM instruction set, and left as future work to formalize the meaning of the graph given to their WCET analyzer and prove that it captures the control paths of the binary \cite{sewell2017high}; the next theorem is such a proof for the graphs of Definition~\ref{def:cycle-certificate}, relative to a specification written for the purpose.

\begin{defn}[Runs of machine code]
\label{def:machine-run}
A \emph{specification} of the processor maps an address $a$ to the instruction stored there, if any, and its successors $\mathit{succ}(a)$: triples $(a', w, G)$ of a next address ($a' = \bot$ for a return), a number of cycles and a list of called functions (by entry address); a conditional branch has two successors, taken and not taken, with their cycles, and an instruction whose effect on the program counter is not modelled has none. A function may also list \emph{jumps through a table}, each an address with a set of targets; the successors of a listed jump are its targets, with the cycles of the jump. A \emph{run} of depth $n + 1$ of the function with entry $e$ is a sequence of steps $(a_i, a'_i, w_i, G_i)$, $0 \le i \le m$, with $a_0 = e$, $a_{i+1} = a'_i$ and $a'_m = \bot$, where each step is a successor of the instruction at $a_i$ or, if there is no instruction at $a_i$ or its effect is not modelled and it is not a listed jump, arbitrary; together with runs of depth $n$ of the called functions; and such that the rows of the function, on the numbers of indices $i$ with $(a_i, a'_i) = (u, v)$ for pairs of addresses $(u, v)$, hold. Its cost is $\sum_i (w_i + c_i)$, with $c_i$ the cost of the runs called at step $i$.
\end{defn}

\begin{thm}[Cycle bounds of machine code]
\label{thm:machine-cycles}
Let each function of a program be given by its entry address, a finite set $V$ of addresses and rows on pairs of addresses, and build its graph: a node per address of $V$ and an exit node, an edge per successor of the instruction at each address of $V$ (per target, for a listed jump), with its cycles and its callees, and the rows carried over to the edges by their pairs of addresses. Suppose that $V$ contains the entry and the successors of its instructions, that every address of $V$ holds an instruction whose effect is modelled or a listed jump, that every callee is the entry of a function of the program, and that the graphs have certificates for the claimed bounds. Then every run of every function costs at most its claimed bound.
\end{thm}

\begin{proof}
By induction on the depth of the run. The addresses of a run stay in $V$, since the entry is in $V$ and $V$ contains the successors of its instructions; each step is therefore a successor of a modelled instruction, hence an edge of the graph with the same cycles, and the steps form a walk from the entry node to the exit node of the same cost, on which each row has the value it has on the run. By the induction hypothesis the calls of each step are runs of the callees' graphs, and Proposition~\ref{prop:cycles} applies.
\end{proof}

The specification is trusted, as an instruction-set semantics is. For the Cortex-M0 it is a decoder of the ARMv6-M Thumb encodings with the cycles of the processor's instruction summary (zero-wait-state memory), written for the purpose (\lean{Thumb.\allowbreak decode}): no data are modelled, so both successors of a conditional branch are allowed, and an instruction that writes the program counter otherwise than by a branch, a call or a return has no successor and makes the check fail, unless it is a listed jump. That a listed jump goes only to its targets is not part of the specification: it is proved from the data semantics (Proposition~\ref{prop:jump-table}), and the checker accepts listed jumps only with the loop certificates of Section~\ref{sec:loops}. A model of the instruction set validated against hardware, such as the HOL4 model of ARMv7 of Fox and Myreen \cite{fox2010trustworthy}, would make that specification more trustworthy; it is not used here. The rows remain an assumption on the executions here; Section~\ref{sec:loops} removes it for the loops that it certifies. Maroneze et al.\ prove the generation of their IPET program from the semantics of CompCert's assembly language \cite{maroneze2014formally}, where the code is the compiler's output rather than a binary.

\section{Loop Bounds on the Machine Code}
\label{sec:loops}

The rows of Definition~\ref{def:cycle-certificate} are assumptions on the runs. In Theorem~\ref{thm:machine-cycles} they come from the source program, through the matching of source and binary loops, and they are the remaining link between the source analysis and the bound of the binary. To check them on the machine code itself we need a semantics with data: the program counter alone allows both successors of every conditional branch, and no loop is bounded in it. We model the effect of each instruction on the low registers $r_0, \dots, r_7$, on the flags, on the stack pointer and on the memory, of which only the words of the function's frame accessed through SP are tracked.

\begin{defn}[Data runs]
\label{def:data-run}
A \emph{state} holds eight words $r_0, \dots, r_7 \in [0, 2^{32})$, four flags $N, Z, C, V$, a stack pointer $\mathit{sp}$ and a memory, a function from addresses to words. The specification gives every instruction an \emph{effect}: for each low register, that it is kept, that it receives $(a + b) \bmod 2^{32}$ or $(a - b) \bmod 2^{32}$ for operands $a, b$ that are registers or constants, that it receives $a \cdot 2^{s}$ if $a \cdot 2^{s} < 2^{32}$ and any value otherwise, for a register $a$ and a constant $s$, or that it receives any value; for the flags, that they are kept, that they are those of the addition $a + b$ or of the subtraction $a - b$ as ARM's AddWithCarry computes them, or that they are any; for the stack pointer, that it is kept, moved by a constant, or any; and for memory, that it is kept, that the word at $\mathit{sp} + i$ receives a register (a store through SP), that only the $n$ words below $\mathit{sp}$ change (\texttt{PUSH}), or, for a store through another register or a call, that any word may change, except, if the function's certificate assumes it (\emph{frame}), the words of the function's frame, between $\mathit{sp}$ and the stack pointer $\mathit{sp}_0$ at the function's entry; a load through SP gives the word at $\mathit{sp} + i$. A conditional branch is \emph{taken} exactly when its condition, a function of the flags, holds, and a jump \texttt{ADD PC, rm} or \texttt{MOV PC, rm} ($m < 8$) goes to $(a + 4 + r_m) \bmod 2^{32}$ or to $r_m$, bit~0 cleared, in 3~cycles; the instructions \texttt{ADD rd, PC}, \texttt{ADR} and the loads \texttt{LDR}, \texttt{LDRH}, \texttt{LDRB} from an address of the code also give their exact values ($r_d + a + 4$, the aligned address of \texttt{ADR}, the word, halfword or byte of the code); the code includes the read-only data that follow it, and neither is written during a run. The result of \texttt{ANDS rd, rm} is at most each operand, that of \texttt{LSRS rd, rm, \#s} at most $2^{32 - s} - 1$ (and $0$ for $s = 0$, a shift by 32), and that of \texttt{UXTB} and \texttt{LDRB} at most $255$. A \emph{data run} of a function is a run of Definition~\ref{def:machine-run} together with a state before each step, such that each step's instruction is modelled, the step is its taken successor if and only if the condition holds (for a conditional branch), and each state is an effect of the previous instruction on the previous state; a call is a single step, whose effect is that of the call instruction. Only the rows designated as trusted need to hold on a data run.
\end{defn}

For the Cortex-M0 the effects are exact for the additions and subtractions of low registers and immediates (\texttt{ADDS}, \texttt{SUBS}), the moves of an immediate or between low registers, the comparisons \texttt{CMP} and \texttt{CMN}, the shift \texttt{LSLS rd, rm, \#s} ($\mathit{rm} \cdot 2^{s}$ when it does not overflow, any value when it does; its flags are any; with $s = 0$ it is a move between low registers), the negation \texttt{NEGS} ($0 - x$, with the flags of that subtraction) and the complement \texttt{MVNS} ($(2^{32} - 1) - x$; its flags are any), the write-back of \texttt{STM} and \texttt{LDM} (the base register advances by four bytes per register of the list, \texttt{STM} writes no register of its list, and \texttt{LDM} gives any value to the registers of its list, including the base if it is in the list), and the load of a word of the literal pool (\texttt{LDR rt, [pc, \#i]}), which gives the word of the code at that address when it lies in the certificate's code; every other instruction writes the low registers it may write, with any value, and the flags, with any value, if it may change them, except for the exact values of Definition~\ref{def:data-run} (\texttt{ADD rd, PC}, \texttt{ADR}, loads from the code). The effect of the literal load, and of the loads from the code, assumes that the code is not written during a run, as the PC-level semantics already does by decoding the instructions from a fixed code. A load through SP (\texttt{LDR rt, [sp, \#i]}) gives the word stored there; a load through another register gives any value. \texttt{SUB}/\texttt{ADD SP}, \texttt{PUSH} and \texttt{POP} move the stack pointer by their constants, and a write of SP from a register makes it any. A call may change $r_0, \dots, r_3$ and the flags; it may change every low register and the stack pointer unless the callees are assumed to preserve $r_4, \dots, r_7$ and SP, as the Arm procedure call standard requires of them (\emph{abi}; Section~\ref{subsec:conv} replaces this assumption by a check of the callees' code). The assumption \emph{frame} is weaker than it may seem only if the function takes no address of its frame: then, in a C program without undefined behaviour, no other code can write the frame (a callee's own frame is below the caller's stack pointer). The checker therefore refuses \emph{frame} for a function that writes an address derived from SP into a register (\texttt{ADD rd, SP, \#i}, \texttt{MOV rd, SP}, \texttt{ADD rd, SP}), with one exception: the frame pointer set by the usual prologue \texttt{ADD r7, SP, \#i}. In a function with that instruction, the checker refuses \emph{frame} if any instruction may read $r_7$ other than as the base address of a load, or than a \texttt{PUSH} at the function's entry (which saves the caller's $r_7$); the registers that an instruction may read are over-approximated from its encoding. The address in $r_7$ then reaches no store and no other register, so that the argument above still applies (a callee receives $r_7$, as it receives SP, and the assumption covers calls in both cases). These checks only restrict the functions for which the assumption may be stated: the soundness proof does not use them. Both assumptions are stated in the certificate and reported by the checker.

\begin{defn}[Loop certificate]
\label{def:loop-certificate}
Let $f$ be a function with entry $e$ and instruction addresses $V$. A \emph{loop certificate} of $f$ gives a list of \emph{tracked words}, distinct offsets $0 < k_1, \dots, k_S < 2^{20}$, multiples of 4 (the words at $\mathit{sp}_0 - k_j$), and every $a \in V$ an offset $\delta_a$ of the stack pointer from $\mathit{sp}_0$ (or none), a \emph{kind} $k_a \in \{\mathrm{none}, \mathrm{add}, \mathrm{sub}\}$, an \emph{invariant} given as a list $I_a^0, \dots, I_a^{n_a}$ of \emph{disjuncts}, finite sets of linear constraints over $r_0, \dots, r_7$, two \emph{ghost} variables $g_1, g_2$ and the tracked words $m_1, \dots, m_S$ (usually $n_a = 0$), a finite set $C_a$ of \emph{congruences} $x \equiv r \pmod{2^{e}}$, $e \le 32$, on the same variables, and, for each case of each transition from $a$ (the transition, its wrap-around case $m$, a fact of its branch and a disjunct $I_a^i$), the index of a disjunct of the next instruction's invariant ($0$ if none is given); and each certified loop $L$ a header $h \in V$, a body $B \subseteq V$, a scale $D \ge 1$, a bound $N \ge 0$ and an affine potential $\Phi_a$ for each $a \in B$ ($\Phi_a = 0$ for $a \notin B$ and for a return). A state $\sigma$ with ghosts $g$ satisfies the certificate at $a$ if $\sigma$, $g$ and the words $m_j$ of $\sigma$'s memory at $\mathit{sp}_0 - k_j$ satisfy one of the disjuncts $I_a^i$ and $C_a$, $\mathit{sp} = \mathit{sp}_0 + \delta_a$ (modulo $2^{32}$) when $\delta_a$ is given, and, when $k_a \ne \mathrm{none}$, the flags of $\sigma$ are those of $g_1 + g_2$ or $g_1 - g_2$ according to $k_a$. The certificate is \emph{valid} if
\begin{enumerate}
  \item $k_e = \mathrm{none}$, $\delta_e = 0$ or none, $C_e = \emptyset$, $I_e^0$ holds in every state, and $\Phi_e \le \beta$ in every state, where $\beta = N$ if $e \in B$ and $\beta = 0$ otherwise;
  \item for every transition from $a$ to $a'$ allowed by the effect and the condition at $a$, from a state that satisfies the certificate at $a$ with ghosts $g$, in its disjunct $I_a^i$: the offsets agree with the effect on the stack pointer ($\delta_{a'}$ is none, or $\delta_a$ is given and $\delta_{a'}$ is $\delta_a$ moved by the effect), and the state after satisfies the certificate at $a'$ with ghosts $g'$, in the disjunct of $a'$ that the certificate names for the case, where $g' = g$ if the flags are kept, $g'$ are the operands if the flags are those of an addition or subtraction (and then $k_{a'}$ is $k_a$, the new kind or none), and $k_{a'} = \mathrm{none}$ if the flags become any; and, for every certified loop,
  \[
    \kappa_L(a, a') + \Phi_{a'}(\text{after}) - \Phi_a(\text{before}) \;\le\; 0, \qquad
    \kappa_L(a, a') = \begin{cases} D & a' = h,\ a \in B,\\ -N & a' = h,\ a \notin B,\\ 0 & \text{otherwise.}\end{cases}
  \]
\end{enumerate}
The \emph{row} of $L$ is $\sum_{u \in B} D\, x_{(u,h)} - \sum_{u \notin B} N\, x_{(u,h)} \le \beta$ on the counts of the transitions $(u, h)$.
\end{defn}

The row says that the back edges of the loop are taken at most $N/D$ times per entry; with $D = 1$ it is the row of Section~\ref{sec:graphs}. Each implication of condition~(2) on the constraints and the potentials is checked by a Farkas certificate, with the constraints of the disjunct $I_a^i$, the range $[0, 2^{32})$ of every variable and the facts of the transition as premises. The congruences need no certificate: the value of a variable after the instruction is an affine expression $\sum_v k_v x_v + c$ of the variables before it (the same expression as for the constraints, with its wrap-around term $m\, 2^{32}$; for a shift that overflows, the new value of the register, a variable without congruence), and the checker accepts $x' \equiv r \pmod{2^e}$ if, for every $v$, $2^e$ divides $k_v 2^{e_v}$, where $x_v \equiv r_v \pmod{2^{e_v}}$ is the congruence of $C_a$ on $v$ ($e_v = 0$ if there is none), and $\sum_v k_v r_v + c \equiv r \pmod{2^e}$. Since $e \le 32$, wrapping modulo $2^{32}$ preserves the congruences. After \texttt{LSLS rd, rm, \#s} it also accepts $r_d' \equiv 0 \pmod{2^e}$ for $e \le s < 32$, overflow or not, the value being $(r_m 2^s) \bmod 2^{32}$. The tracked words after an instruction are computed by the checker from the offsets: a store through SP at offset $\delta_a + i$ writes the word $j$ with $k_j = -(\delta_a + i)$ and no other (the offsets are small, so distinct offsets are distinct addresses modulo $2^{32}$); \texttt{PUSH} may write the words in $[\delta_a - 4n, \delta_a)$; a store through another register or a call keeps the words of the frame ($k_j \le -\delta_a$) under the assumption \emph{frame} and may write every word otherwise; a load through SP at offset $-k_j$ gives $m_j$. The facts of a sum or difference written into a register are a case split on its wrap-around: the value is $a + b + m \cdot 2^{32}$ for some $m \in \{-1, 0, 1\}$ with $0 \le a + b + m\, 2^{32} < 2^{32}$, and the certificate proves the goal in each case (the cases that the invariant excludes by a contradiction). The facts of a shift $a \cdot 2^{s}$ are a case split on overflow: either $a \cdot 2^{s} \le 2^{32} - 1$ and the value is $a \cdot 2^{s}$, or $a \cdot 2^{s} \ge 2^{32}$ and the value is any word. The facts of a conditional branch are a disjunction of linear constraints on the ghosts, by the following lemma.

\begin{lem}[Conditions as linear facts]
\label{lem:cond-facts}
For words $g_1, g_2$ and each of the conditions EQ, NE, CS, CC, HI, LS, GE, LT, GT, LE, whether it holds on the flags of $g_1 - g_2$ is a disjunction of conjunctions of linear constraints on $g_1, g_2$; for EQ, NE, CS, CC, HI and LS, so is whether it holds on the flags of $g_1 + g_2$. For instance GE holds on the flags of $g_1 - g_2$ if and only if $g_1, g_2 < 2^{31} \wedge g_1 \ge g_2$, or $g_1, g_2 \ge 2^{31} \wedge g_1 \ge g_2$, or $g_1 < 2^{31} \le g_2$; and NE holds on the flags of $g_1 + g_2$ if and only if $1 \le g_1 + g_2 \le 2^{32} - 1$ or $g_1 + g_2 \ge 2^{32} + 1$.
\end{lem}

\begin{proof}
From the definition of AddWithCarry: for the subtraction, $C$ holds if and only if $g_1 \ge g_2$, $Z$ if and only if $g_1 = g_2$, and $N = V$ if and only if $\mathrm{sgn}(g_1) \ge \mathrm{sgn}(g_2)$, where $\mathrm{sgn}$ reads a word in two's complement; the signed comparison is then split by the signs. For the addition, $C$ holds if and only if $g_1 + g_2 \ge 2^{32}$ and $Z$ if and only if $g_1 + g_2 \in \{0, 2^{32}\}$. The other conditions (MI, PL, VS, VC) give no fact.
\end{proof}

A counter stepped by $4$ that leaves its loop when it equals $40$ satisfies $g_1 \le 40$ and NE, that is $g_1 \le 39$ or $g_1 \ge 41$, at the back edge; no linear invariant follows that bounds it by $36$. Congruences in the sense of Granger \cite{granger1989static}, restricted to powers of two, give the missing step: if $g_1 \equiv g_2 \pmod{2^e}$, then $g_1 - g_2$ is a multiple of $2^e$, and every fact $\pm(g_1 - g_2) + c \le 0$ of Lemma~\ref{lem:cond-facts} may be replaced by $\pm(g_1 - g_2) + 2^e \lceil c / 2^e \rceil \le 0$ (the reduction by which Granger combines congruences with intervals, rounding each bound to the class \cite[Proposition~6.1]{granger1989static}, here applied to the difference $g_1 - g_2$). The checker computes the largest such $e \le 32$ from the congruences of $C_a$ on the ghosts and rounds the facts of the branch before they enter the premises. The facts of the flags of an addition are on the sum $g_1 + g_2$; when the congruences make the sum a multiple of $2^e$, its facts $\pm(g_1 + g_2) + c \le 0$ are rounded in the same way. Clang compiles some loops over arrays with a byte offset counted from $2^{32} - 16$ up by $4$ to zero (\texttt{ADDS r3, \#4}, \texttt{BNE}): in the case of the wrap-around, the loop continues only if $g_1 + 4 \ge 2^{32} + 1$, rounded to $g_1 + 4 \ge 2^{32} + 4$, which the bound $g_1 \le 2^{32} - 4$ contradicts.

\begin{thm}[Loop bounds of machine code]
\label{thm:machine-loops}
Suppose that the functions of a program satisfy the hypotheses of Theorem~\ref{thm:machine-cycles}, with, as rows, the rows designated as trusted and the rows of the certified loops of valid loop certificates. Then every data run of every function on which the trusted rows hold costs at most its claimed bound. In particular, if no row is trusted, every data run does.
\end{thm}

\begin{proof}
Let $(a_i, \sigma_i)_{0 \le i \le m}$ be the steps and states of a data run of $f$. By condition~(1) the certificate holds at $a_0 = e$ with any ghosts, and by condition~(2) and induction on $i$ it holds at every $a_i$, with ghosts and disjuncts chosen along the run (at $a_{i+1}$, the disjunct that the certificate names for the case of the transition, which exists since the state after the transition lies in it); condition~(2) also gives, for every certified loop $L$, $\kappa_L(a_i, a_{i+1}) \le \Phi_{a_i}(\sigma_i) - \Phi_{a_{i+1}}(\sigma_{i+1})$, with $\Phi = 0$ after the return. Summing, the potentials telescope:
\[
  \sum_{i} \kappa_L(a_i, a_{i+1}) \;\le\; \Phi_e(\sigma_0) \;\le\; \beta,
\]
and the left-hand side is the value of the row of $L$ on the counts of the run. The addresses of the tracked words are fixed by $\mathit{sp}_0$, the stack pointer of the first state, so that a word keeps its meaning along the run; the offsets give the stack pointer at every step, and the checker's computation of the tracked words after each instruction is correct because the addresses $\mathit{sp}_0 - k_j$ and $\mathit{sp}_0 + \delta_a + i$ coincide modulo $2^{32}$ exactly when $k_j = -(\delta_a + i)$, all these offsets being below $2^{21}$ in absolute value. Hence every row holds, the data run is a run of Definition~\ref{def:machine-run}, and Theorem~\ref{thm:machine-cycles} applies; the calls are data runs of smaller depth, by induction on the depth.
\end{proof}

\paragraph{Jumps through a table.} Clang compiles a \texttt{switch} on dense cases for the Cortex-M0 into a range check and an indexed jump, for instance \texttt{CMP r0, \#3; BHI} default; \texttt{ADD r0, PC; LDRB r0, [r0, \#4]; LSLS r0, r0, \#1; ADD PC, r0}, the table of byte offsets following the jump, or \texttt{LSLS r0, r5, \#2; ADR r1, \#4; LDR r0, [r1, r0]; MOV PC, r0} with a table of addresses. The checker establishes the targets of such a jump by evaluating its window on the code. The table may also lie in the read-only data after the code, as in libgcc's divisions (\texttt{LDR r1, =table; LSLS r3, r3, \#2; LDR r3, [r1, r3]; MOV PC, r3}); the certificate then holds these data with the code. When the compiler knows the range of the index, it emits no range check: \texttt{switch (x \& 7)} with all eight cases becomes \texttt{MOVS r5, \#7; MOV r3, r2; ANDS r3, r5; ADD r3, PC; \dots}; the bound then comes from the instruction that computes the index (second part of the proposition). When the index is bounded only by a loop, as in \texttt{bitcount}, whose \texttt{switch (i)} sits in a loop on $i < 8$ (\texttt{LSLS r0, r6, \#2} stores $4i$ in the frame, reloaded before \texttt{ADR r0; LDR r0, [r0, r1]; MOV PC, r0}), or by constants set on several paths, as in \texttt{sha}'s second \texttt{switch}, the bound comes from the certified invariant at an instruction of the window (third part).

\begin{prop}[Jumps through a table]
\label{prop:jump-table}
Let a function list a jump at address $p$ with targets $T$ and a window starting at $s$, $s + 4 \le p$, such that the instruction at $s$ writes no register and sets the flags of $r_I - n$ for a low register $r_I$ and a constant $n$, the instruction at $s + 2$ writes no register and is either a \texttt{BHI} whose target is not $s + 4$ (the window then continues at $b = s + 4$) or a \texttt{BLS} to $s + 6$, over a branch to the default case at $s + 4$ (the window continues at $b = s + 6$), and the instruction at $p$ is \texttt{ADD PC, rm} or \texttt{MOV PC, rm} with $m < 8$. Suppose that (1) for every $x \le n$, the instructions from $b$ to $p - 2$, evaluated in this order on partial registers (each a word or unknown), starting from $r_I = x$ and all other registers unknown, by their effects and exact values (an unknown operand gives an unknown result), give a known value $v$ to $r_m$, and the address the jump computes from $v$ is in $T$; and (2) no transition of the function leads to $s + 2$ except from $s$, to $b$ except from $s + 2$, or to an address of $b + 2, \dots, p$ except from the address before it, the entry is none of them, and no listed jump is at $s + 2$. Then in every data run of the function, every step at $p$ goes to an address of $T$; hence every step of a data run follows an edge of the function's graph, and Theorems~\ref{thm:machine-cycles} and \ref{thm:machine-loops} hold with the listed jumps.

The same conclusion holds for a jump listed with a window from $s$ to $p$ bounded at an address $q$, $s \le q < p$, if (1') the instructions from $s$ to $q - 2$, evaluated from all registers unknown, give the operand of a bound of the instruction at $q$ a known value (an \texttt{ANDS} whose other operand is known, or an \texttt{LSRS}, \texttt{UXTB} or \texttt{LDRB}, Definition~\ref{def:data-run}), so that the register $r_d$ that it writes is at most a constant $B$ (the checker requires $B < 1024$), and for every $x \le B$, the instructions from $q + 2$ to $p - 2$, evaluated from the registers after $q$ with $r_d = x$, give the jump an address in $T$; and (2') no transition leads to an address of $s + 2, \dots, p$ except from the address before it, and the entry is none of them.

It holds also for a jump listed with a window from $s$ to $p$, an address $q$, $s \le q \le p$, a low register $r_d$ and bounds $\ell \le u < 1024$, if the function has a valid loop certificate (Definition~\ref{def:loop-certificate}), (2') holds, and (1'') every disjunct $I_q^i$ of the invariant at $q$ contains a constraint $r_d - u' \le 0$ with $u' \le u$ and, if $\ell > 0$, a constraint $\ell' - r_d \le 0$ with $\ell' \ge \ell$, and for every $x$ with $\ell \le x \le u$ and $x \equiv r \pmod{2^e}$, where $r_d \equiv r \pmod{2^e}$ is the congruence of $C_q$ on $r_d$ ($e = 0$ if there is none), the instructions from $q$ to $p - 2$, evaluated from the registers that those from $s$ to $q - 2$ give from all registers unknown, with $r_d = x$, give the jump an address in $T$.
\end{prop}

\begin{proof}
By strong induction on the position of the steps in the run: every step follows an edge of the graph. A step at $p$ is preceded, by~(2) applied to the steps before it, which follow edges by the induction hypothesis, by steps at $p - 2, \dots, b, s + 2, s$, the entry not being in the window. The instruction at $s$ sets the flags of $r_I - n$ and keeps the registers. The \texttt{BHI} is not taken, since its step goes to $s + 4$ and not to its target; the \texttt{BLS} is taken, since its step goes to $s + 6$ and not to $s + 4$ (no listed jump at $s + 2$ adds other transitions). In both cases $r_I \le n$ by the flags of the subtraction, and $r_I$ is unchanged at $b$. Along the window, every known value of the evaluation from $x = r_I$ is the value of its register in the run, by the effects and the exact values of Definition~\ref{def:data-run}; at $p$, $r_m = v$, so the jump goes to the address computed from $v$, which is in $T$ by~(1). For the second form, (2') gives steps at $p - 2, \dots, s$; the evaluation from $s$ is exact up to $q$, so the bound of $r_d$ after $q$ is the constant $B$, and the value $x = r_d$ in the run is one of those evaluated by~(1'). For the third form, the induction carries the certificate along with the edges: a step at $p$ is preceded by steps at $p - 2, \dots, s$, by~(2'), and in particular by a step at $q$; the steps before it follow edges by the induction hypothesis, so that condition~(1) of Definition~\ref{def:loop-certificate} at the entry and condition~(2) along these steps give, by induction on their number, that the state at $q$ satisfies the certificate at $q$: some disjunct $I_q^i$ and the congruences $C_q$ hold, so $\ell \le r_d \le u$ and $r_d \equiv r \pmod{2^e}$, and $x = r_d$ is one of the values evaluated by~(1''). Condition~(2) is used only on steps that follow edges, which is all that a valid certificate provides.
\end{proof}

The evaluation does not depend on the instructions of the window beyond their effects: a call or a store between the branch and the jump only makes registers unknown. A table whose index no range check, no instruction of the window and no certified invariant bounds within 1024 values is refused, and so is a table in writable data or in read-only data that do not follow the code. The producer meets a circularity in the third form: the graph needs the targets, the invariant the graph. It builds the graph with the targets found so far (none at first), computes the zones and congruences at $q$, evaluates the window on that range and repeats until the targets no longer change; the checker sees only the result, and the certificate is valid for the graph with these targets. The calling convention (Section~\ref{subsec:conv}) is checked along the same transitions, the listed targets included.

The proof uses neither the order of the instructions nor the structure of the loops: a certified loop need not be a natural loop, and its body need not be the set of instructions that the producer found. Only the generated row matters. In the checker, the rows of the certified loops are generated from the certificate, after the trusted ones, and the graph of Theorem~\ref{thm:machine-cycles} is built with them, so that the cycle certificate is checked against the rows that Theorem~\ref{thm:machine-loops} proves.

The certificates are produced by an analysis of the binary, untrusted (script \lean{benchmarks/cycles/binloops.py}): the kinds by a forward data-flow analysis, the offsets of the stack pointer by a forward data-flow analysis (an offset is dropped where two paths disagree), the tracked words as the words accessed through SP inside loops, the congruences by a forward data-flow analysis on the checker's own post-expressions (moduli $2^e$, joined by lowering $e$ until the residues agree; the lattice has finite height), the invariants by an analysis of differences of registers and tracked words (a zone domain, with widening at loop headers and descending iterations, in which the facts of each branch are rounded with the congruences), restricted to the variables that the potentials and the loop tests involve, and the potentials by trying $\pm x$ for every register and tracked word and $\pm (x - y)$ for the variables compared at the loop's exits, with per-instruction constants computed as longest paths and the scale $D$ chosen among $1, 2, 4, 8$. When no such potential is found, or its bound exceeds the source bound, a linear program searches for a potential with its own linear part at each instruction: its unknowns are the coefficients of the potentials, the bound and the Farkas multipliers of every obligation of condition~(2), whose premises (the invariants, the ranges and the facts of the transition) are fixed, so that the obligations are linear in the unknowns, and it minimizes the bound. This is the synthesis of linear ranking functions from linear constraints by duality \cite{colon2001synthesis}, with an affine function per control point as in Alias et al.\ \cite{alias2010multi}; the rational solution is rounded, multiplied by the least common denominator of its coefficients, which becomes the scale $D$, and proved again in exact arithmetic (the linear-programming solver, HiGHS through scipy, is limited to 5~seconds per program, since the constants $2^{32}$ of the ranges make these programs badly scaled). It is needed when the counter moves between a register and a word of the frame inside the loop: no single linear part decreases along such a loop, since the counter is in the register at some instructions and in the word at others. An addition or subtraction of a constant that may wrap around splits the zone in two, the case without and the case with wrap-around, which are carried separately along the straight line of instructions up to the next conditional branch and joined after it (trace partitioning, with the partitions merged at a later control point \cite{mauborgne2005trace}); the instructions in between get one disjunct per case, and each case of a transition is sent to the disjunct of the same case. The run-time division routines need it: they count \texttt{for (i = 31; i >= 0; i--)}, compiled as \texttt{SUBS r3, r3, \#1}, \texttt{ADDS r4, r3, \#1}, \texttt{BHS} to the exit, which is taken only once $r_3$ has wrapped around to $2^{32} - 1$; between the subtraction and the branch, $r_3 \le 30$ or $r_3 = 2^{32} - 1$, which no conjunction of linear constraints expresses (for one variable it is a wrapped interval \cite{gange2015interval}; here the disjuncts also relate $r_3$ to its copy in the operand of the flags and to $r_4$). The potential then follows the counter through its copies: $r_3 + 1$, the old value of $r_3$ (the ghost $g_1$) after the subtraction, $r_4 + 1$ at the branch, in both disjuncts ($r_4 = 0$ after the wrap-around). HiGHS often declares these linear programs infeasible, or fails, because their constants span nine orders of magnitude; the producer then retries without HiGHS's presolve, then with each constant $c$ written $q\, 2^{31} + r$, $|r| \le 2^{30}$, and the two parts of every constant row required separately (a stronger, well-scaled program), and finally with the linear parts fixed to a variable that follows a bounded variable of the header through the body (a variable whose value after each instruction is the value before it plus a constant, in every disjunct and wrap-around case), only the constants being unknown; these constants are computed first as longest paths over the transitions of each disjunct, as for the potentials $\pm x$, and only then by the linear program. Widening uses as thresholds the constants of the code, their neighbours and, for each addition of a constant $k > 0$ that may wrap around, the values $2^{32} - k - 1$, $2^{32} - k$ and $2^{32} - k + 1$ at which it wraps, so that a counter that leaves its loop when it wraps to zero keeps its bound below $2^{32} - k$. The Farkas multipliers come from shortest paths in the graph of the difference constraints, extended by the bound of one variable that a premise on a sum $x + y$ (the facts of the flags of an addition) gives with a bound of the other, or from a negative cycle when the premises are contradictory, or from a premise that the one-variable bounds contradict (exact, where the linear program loses the difference of two constants near $2^{32}$), or else from a linear program; a constraint or a loop whose obligations cannot be proved is dropped, and a dropped loop keeps its row from the source, trusted.

\subsection{The calling convention, checked}
\label{subsec:conv}

The assumption \emph{abi} is a property of the code of the callees, which the checker has: it can be checked instead of assumed. The Arm procedure call standard requires a function to preserve $r_4, \dots, r_7$ and SP. Compilers do it by saving, with one \texttt{PUSH} at the entry, the registers of $r_4, \dots, r_7$ that the function writes, together with the return address, and by restoring them with a \texttt{POP} of the same list when the function returns. That this restores the registers depends on the saved words not being overwritten in between, which no analysis of the function's own instructions can show for stores through pointers; we state it as a hypothesis on the program.

\begin{defn}[Linked data runs]
\label{def:linked-run}
A \emph{linked data run} of a function from a state $\sigma_i$ to a state $\sigma_o$ is a data run (Definition~\ref{def:data-run}, without \emph{abi}) whose first state is $\sigma_i$ and whose return instruction leads to $\sigma_o$, such that moreover:
\begin{enumerate}
    \item at a call, the callee has a linked data run from the caller's state to a state $\sigma$, and the caller continues with the low registers and the stack pointer of $\sigma$;
    \item \texttt{PUSH} stores the registers of its list, in ascending order, in the words below the stack pointer, and \texttt{POP} loads them back from the words above it (ARMv6-M);
    \item \emph{no stack smashing}: the words in which the function saves registers at its entry are changed neither by a store through a register other than SP nor by a call.
\end{enumerate}
\end{defn}

Hypothesis~(3) holds for every program that writes memory only within the objects it addresses; a write that breaks it would also overwrite the return address saved with the registers, which is what stack protectors detect. It is implied by \emph{frame} for the functions to which \emph{frame} applies, and, unlike \emph{abi}, it does not depend on how the callees were compiled.

\begin{defn}[Calling convention]
\label{def:convention}
Let $f$ be a function with entry $e$ whose loop certificate gives the offsets $\delta_a$ of the stack pointer, and let $P$ and $n$ be the low registers and the number of words pushed by the instruction at $e$ if it is a \texttt{PUSH} ($P = \emptyset$ and $n = 0$ otherwise). The function \emph{keeps the convention} if $\delta_e = 0$, no transition leads to $e$ (the targets of the listed jumps through tables included), and every other instruction $a$ satisfies: if it returns, it is no listed jump and it is a \texttt{POP} of $P$ and the program counter with $\delta_a = -4n$, the entry having pushed the link register, or a \texttt{BX} with $n = 0$ and $\delta_a = 0$; otherwise it writes no register of $\{r_4, \dots, r_7\} \setminus P$, it stores through SP at an offset $\delta_a + i < -4n$ if it stores through SP, and $\delta_a \le -4n$ if it is a \texttt{PUSH}.
\end{defn}

\begin{thm}[The calling convention, checked]
\label{thm:machine-conv}
Suppose that the loop certificates of the functions of a program are valid with \emph{abi} and that every function keeps the convention. Then every linked data run of a function from $\sigma_i$ to $\sigma_o$ is a data run with \emph{abi}, and $\sigma_o$ has the registers $r_4, \dots, r_7$ and the stack pointer of $\sigma_i$. Consequently, the conclusion of Theorem~\ref{thm:machine-loops} holds for every linked data run, without the assumption \emph{abi}.
\end{thm}

\begin{proof}
By induction on the depth of calls. The calls of a run are linked runs of smaller depth, which, by the induction hypothesis, restore $r_4, \dots, r_7$ and the stack pointer; every call step therefore has the effect of a call with \emph{abi}, and the run is a data run with \emph{abi}. By Proposition~\ref{prop:jump-table}, whose argument uses only the effects with \emph{abi} just established, every step of the run follows a transition of the function, a listed target at a jump through a table. For the second claim, we show along these transitions that (i) the stack pointer is $\mathit{sp}_0 + \delta_a$ whenever $\delta_a$ is given, by the validity of the offsets; (ii) the registers of $\{r_4, \dots, r_7\} \setminus P$ keep their values at the entry, since no instruction writes them and calls preserve them; and (iii) after the entry, the word at $\mathit{sp}_0 - 4n + 4\rho_k$, where $\rho_k$ is the position of $r_k$ in the list $P$, holds the value of $r_k$ at the entry, for every $r_k \in P \cap \{r_4, \dots, r_7\}$. The \texttt{PUSH} at the entry establishes (iii) by~(2). A store through SP at an offset below $-4n$, and a \texttt{PUSH} with the stack pointer at or below $\mathit{sp}_0 - 4n$, write other words, since all offsets are below $2^{21}$ in absolute value and distinct offsets are distinct addresses modulo $2^{32}$; stores through other registers and calls keep these words by~(3). At a \texttt{POP} of $P$ with the stack pointer at $\mathit{sp}_0 - 4n$, (2) loads every register of $P$ from its saved word, the other registers keep their values by~(ii), and the stack pointer becomes $\mathit{sp}_0$; at a \texttt{BX}, nothing has been saved and the stack pointer is $\mathit{sp}_0$.
\end{proof}

The check needs the offsets of the stack pointer at the returns and at the stores through SP, which the producer now gives for every function, with loops or not. The checker program applies it when a certificate states \emph{abi}: if every function of the program keeps the convention, it reports the preservation of $r_4, \dots, r_7$ and SP as checked, under hypothesis~(3); otherwise it reports \emph{abi} as an assumption and names the functions that do not keep the convention.

\section{Evaluation}
\label{sec:evaluation}

The binaries are those of the 29 kernels of TACLeBench, as in Paper~II \cite[Section~\ref{II-subsec:tacle}]{crosscheckTool}, whose loop bounds in the calling contexts of the programs are the source bounds used here.

\paragraph{Cycles on a predictable processor.}
On a processor whose instruction timings do not depend on the execution history, cost bounds can become time bounds. For the Arm Cortex-M0 (in order, without cache or branch prediction; zero-wait-state memory assumed), the kernels were compiled with Clang at \texttt{-O1} without loop unrolling, and each function of the binary was bounded in cycles by the timing schema applied to its control-flow graph: a loop costs its iteration bound times its longest iteration plus its longest exit path, a call its callee's bound, a conditional branch the taken case. The iteration bound of a binary loop is the in-context bound that CrossCheck derived for the source loop it is matched with through the debug line table: of the lines of the binary loop's own instructions (those outside its nested loops), the ones that lie in some source loop, from the loop statement's first line to its last, are collected, and the innermost source loop that holds all of them is chosen; lines outside every loop, which the compiler gives for instance to code it moved, are ignored, and if no source loop holds all of them, the largest bound among the loops whose first line occurs is taken. The instruction timings are those of the instruction summary of the processor's reference manual; a binary loop is assumed to iterate at most as often as its source loop. Both assumptions were checked on executions: an interpreter of the instruction set, written for this purpose and tested against the C semantics, runs each kernel (each program checks its own result), counts cycles with the same table, and compares every call of a bounded function and every binary loop with its bound. Integer division, which the M0 lacks, is compiled from a small C implementation analysed like the kernels; floating point, done in software on the M0, is not covered in these figures (its routines have no code in the binary; see the paragraph ``Floating point on its machine code''). Of the 267 functions of the 29 kernels, 107 receive a cycle bound, including the entry functions of 6 kernels (\texttt{binarysearch}, \texttt{bsort}, \texttt{countnegative}, \texttt{isqrt}, \texttt{jfdctint}, \texttt{matrix1}), whose bounds are 1.24 to 2.46 times the observed cycles; no bound is exceeded on the 11 calls and 45 loops compared. One of these calls, of a function of \texttt{fft} bounded since logarithmic potentials were added, has a bound 970 times its observed cycles: its three nested loops are bounded by the product of their bounds, while at each outer iteration the count of one inner loop doubles and that of the other halves. Simpler rules failed: the largest bound among all the lines of a loop gave \texttt{matrix1} a bound 105 times the observed cycles, because an instruction outside the loop body inherited the line of another loop, and the line of the back-edge branches matched a loop rotated by GCC with its inner loop (paragraph ``Another compiler''). \texttt{fft} fails its own check on the M0: it converts negative floating-point values to \texttt{unsigned long}, which is undefined in C. The comparison checks the path analysis and the loop matching, not the timing table, which the interpreter shares; the table is compared separately, row by row, with the instruction summary of the processor's reference manual \cite{armCortexM0trm} (a test of 87 rows; instructions without a fixed cycle count there, such as \texttt{WFI}, stop the analysis of their function).

\paragraph{Certified cycle bounds.}
The bound of a binary function can also be checked by the verified checker instead of being trusted from the analysis that computed it. The certificate holds the machine code of the binary, and for each function its entry, the addresses of its instructions, its loop rows (``the back edges of the loop matched with a source loop of bound $N$ are taken at most $N$ times its entry edges'', on pairs of instruction addresses) and the dual solution of the IPET program \cite{li1995ilp}: a linear-programming solver (HiGHS, through scipy; untrusted) gives the multipliers of the rows, which are made integer, and the potentials are computed exactly as longest paths (script \lean{benchmarks/cycles/ipet.py}, option \texttt{--certify} of \lean{run\_cycles.py}). The checker decodes the code itself with a decoder and cycle table written in Lean, builds the control-flow graph of each function (an edge per successor of each instruction; a conditional branch costs 3 cycles taken and 1 not taken), and checks the certificate on it (Theorem~\ref{thm:machine-cycles}). All 107 functions with a cycle bound receive a certificate, and the checker accepts all of them (with the run-time division routines they call, 223 functions, 8{,}333 instructions, checked in 0.15~s); the certified bound is below the timing schema's for 60 functions and equal for 47 (median ratio 0.998, smallest 0.72), never above, and none of the 11 calls compared exceeds it. The entry functions' bounds become 1.06 to 1.92 times the observed cycles. The Lean decoder is compared with the decoder of our interpreter on all $2^{16}$ halfwords (each followed by two second halfwords, for \texttt{BL}) and with llvm-objdump on the 24{,}285 instructions of the kernels' objects, without difference; the comparison with our interpreter covers the cycle table, which is itself compared with the processor's manual (paragraph above). What the check removes from the trusted base is the analysis of the binary, the construction of the graph and the computation of the bound; the specification of the processor (decoder, cycle table, zero-wait-state memory), the copy of the code into the certificate and the loop bounds remain trusted. Maroneze et al.\ prove the generation of their IPET program against the semantics of CompCert's assembly language \cite{maroneze2014formally}, and Sewell et al.\ the decompilation of the seL4 binary against a model of the ARM instruction set validated on hardware \cite{sewell2017high,fox2010trustworthy}; our decoder is not validated on hardware. The run of all kernels with certificates takes 2.8~minutes.


\paragraph{A WCET analyzer on the same binaries.}
The loop bounds were also given to an existing WCET analyzer, OTAWA \cite{ballabriga2010otawa} (version 2.3, built from its public sources), on the same Cortex-M0 binaries, linked into executables. The flow facts that the plugin exports designate loops by source line; on code compiled at \texttt{-O1}, with inlined and rotated loops, OTAWA could not place them, so for this experiment each bound was attached to the header of the binary loop matched with the source loop, as above, designated by function and offset, with the header bounded by one more than the bound of the body. With these flow facts alone, 199 in all, OTAWA computes by IPET a bound for each of the 107 functions that have a cycle bound above, including the 6 entry functions; with one unit per instruction (OTAWA's script \texttt{trivial}), the result bounds executed instructions: of the 11 calls of these functions that the emulated runs complete, none executes more instructions than OTAWA's bound (geometric mean of bound over executed 2.32, the largest ratio, 756, for the same function of \texttt{fft} as above). For cycles, a small program compiled against OTAWA (\lean{benchmarks/cycles/otawa/m0time.cpp}) gives each basic block of OTAWA's control-flow graphs the Cortex-M0 cycles of its instructions, from the timing table above, a conditional branch counted as taken; the Cortex-M0 has no cache and fixed instruction times, so the time of a block is that sum and no pipeline description is needed. OTAWA then bounds the same 107 functions in cycles; none of the 11 calls takes more cycles than its bound (geometric mean 2.69, largest ratio 1{,}054, again in \texttt{fft}). Its bound is never below that of our timing schema: equal for 52 functions and larger for 55 (at most by a factor 1.29, the median ratio being 1.007), the difference coming from the extra execution given to loop headers. The two computations share the timing table and the loop bounds, not the control-flow graphs nor the computation of the bound (IPET in OTAWA, the timing schema here). This checks that the bounds are usable by such a tool and that the binary loop matching is consistent with it; it is not an evaluation of OTAWA.

\paragraph{Other processors.}
For processors whose timing is not modelled here, the loop bounds are given to OTAWA on its own control-flow graphs: OTAWA lists the loop headers and basic blocks of each function, the natural loops of these graphs are matched with the source loops through the debug lines, by the rules above, and written as flow facts (script \lean{benchmarks/cycles/otawa\_flow.py}); a call graph read from the disassembly refuses, for every function that reaches them, recursion, indirect calls and jumps, jump tables, branches into another function and calls of undefined symbols. On the Cortex-M0 binaries above, this matching gives OTAWA the same bounds as the matching of the timing schema for all 107 functions, with no bound below an emulated run (11 calls). The kernels were then compiled for a Cortex-M4 with its single-precision floating-point unit, linked at the addresses of an STM32 microcontroller, and bounded by the two models of OTAWA's \texttt{stm32} script: the STM32F427 (Cortex-M4: pipeline, no cache) and the STM32H753 (Cortex-M7: pipeline, caches, branch history). OTAWA computes a bound for 169 of the 267 functions on the first (8 entry functions) and 163 on the second (7; six analyses do not finish within 300~s); the 107 functions bounded on the Cortex-M0 are all bounded on the Cortex-M4, and 62 functions are bounded only there, 57 of which call a single-precision floating-point function of the run-time library on the Cortex-M0, an instruction of the M4. The other functions reach a loop without a bound (59), an undefined function (24: double precision, \texttt{memcpy}), a loop that matches no source loop (7), recursion (6) or a jump table (2). These bounds are not checked against executions (there is no simulator of these processors here), and the timing models are OTAWA's: its STM32F427 memory description gives every bank a latency of 50 cycles and notes that the values are not from the datasheet, and the Cortex-M4 bounds of the entry functions are 1.4 to 23 times those of the Cortex-M0. The experiment shows that the same source bounds reach another instruction set and another timing analysis through flow facts, not how fast these processors are. The run takes 29 minutes, most of it OTAWA's cache analysis for the Cortex-M7.

\paragraph{Another compiler.} The kernels were also compiled by GCC (arm-none-eabi-gcc 13.2.1, \texttt{-O1} without loop unrolling, jump tables or the replacement of loops by library calls) and analysed in the same way (option \texttt{--gcc} of \lean{run\_cycles.py}). GCC inlines less than Clang, so that the emulated runs compare 73 calls of bounded functions instead of 11, and 65 binary loops. Again 107 functions receive a cycle bound, including the same 6 entry functions (1.31 to 2.51 times the observed cycles), and no bound is exceeded; 112 functions receive a certified cycle bound, and 75 of the 215 loops of the kernels' code and the 29 loops of the division routines are certified on the machine code with both assumptions (66 and 40 of the kernels' loops with \emph{abi} only and with none); why fewer of GCC's loops are certified was not analysed. These binaries revealed two errors in the trusted link between the source and the binary, both found as emulated runs above a bound and both corrected before the figures of this section were computed. The source had been analysed for the data model of the host (64-bit \texttt{long}), and a loop of \texttt{sha} over the bytes of a \texttt{long} received half its iterations on the 32-bit binary: GCC's out-of-line function took 367 cycles against a certified bound of 213 (in Clang's binary the function is inlined and its call was not compared); the source is now analysed for the target. GCC rotates the outer loop of a function of \texttt{fft} so that its back-edge branches carry the line of the inner loop's test, and the earlier rule, which used the line of the back-edge branches first, bounded that loop by the inner loop's bound (11 instead of 1023); the matching by line ranges above replaced it. A third defect was in the producer: on a loop of \texttt{bitcount}, widening followed by the closure of the zone kept raising one bound; after 200 visits of a loop head, the producer now forgets the entries that still grow. That none of these defects appeared on Clang's binaries shows the limits of the check by execution: it covers only the calls and loops that the runs reach out of line.

\paragraph{Stack on the same binaries.}
The stack that a function needs is a high-water mark: the rule over calls is a maximum, $\max(\mathrm{own}(f),\ d + \mathrm{stack}(g))$ over the call sites of $f$, where the rule for cost is a sum \cite[Proposition~\ref{I-prop:stack}]{crosscheckTheory}. On the same Cortex-M0 binaries, the depth of the stack pointer below its value at a function's entry is computed on the control-flow graph (pushes, pops and adjustments of the stack pointer, including its restoration from a register known on every path to hold its entry value minus a constant, as the frame pointer does, a register being kept across a call only if the callee's code saves and restores it or does not write it; the depth must have a single value at every instruction and be $0$ at every return), and the rule is applied bottom-up on the call graph. No loop bound is involved, so more functions are covered than by the cycle bounds: 151 of the 267 functions, including the entry functions of 12 kernels. The interpreter records, for every call, the stack pointer at entry and its lowest value before the return: no bound is exceeded on the 29 calls compared, and 24 bounds equal the observed stack (11 of the 12 entry functions: between 24 and 8{,}648 bytes; \texttt{md5}: 584 bytes against 552 observed). The functions without a bound call software floating-point routines that have no code in this setting, directly or through other functions (107; linked with a run-time library, 256 functions are bounded, paragraph ``Floating point on its machine code''), or are recursive in the binary or call a recursive function (9). Exception frames are not counted.

\paragraph{Time budgets.}
The cost bounds are in abstract units, while industrial budgets are durations. A duration is a property of the machine code on a processor, so it cannot be checked at compile time: an annotation \texttt{CROSSCHECK\_DEADLINE\_US(n)} (or \texttt{\_CYCLES}) is only recorded in the report, together with the loop bounds, and a script checks it after compilation: it computes the cycle bound of the function's code from the plugin's loop bounds, matched with the binary's loops through the debug line table, and a timing model of the processor, converts it with the clock frequency, and fails, so that it can be a post-build step, if the budget is exceeded or the function has no cycle bound. The binary may be the linked program of several translation units, one report per unit: a call is costed from the callee's code in the binary, with the loop bounds of the callee's own report. What it trusts is stated with the result: the timing model, the absence of interrupts during the function, that the compiler preserves the source loops' iteration counts, and the plugin's loop bounds. Two kinds of timing model are supported: the Cortex-M0 model of Section~\ref{sec:evaluation}, and OTAWA \cite{ballabriga2010otawa} for the processors it models, the STM32F427 (Cortex-M4) and STM32H753 (Cortex-M7) of its \texttt{stm32} script, run in a container without network; for OTAWA, the loops of its own control-flow graphs are matched with the source loops, the bounds are given to it as flow facts, and a function that reaches recursion, an indirect call or jump, a jump table or an undefined symbol gets no bound (paragraph ``Other processors'' in Section~\ref{sec:evaluation}). On the Cortex-M0 the cycle bound can also be a certified one, accepted by the verified checker (option \texttt{--certify}; paragraph ``Certified cycle bounds'').

\paragraph{Loop bounds on the machine code.} With option \texttt{--certify-loops}, the loop bounds are also certified on the machine code where the producer finds a certificate (Section~\ref{sec:loops}). The kernels' own functions have 403 natural loops in the binary, and the run-time division routines 116 (four routines per kernel, which shift and subtract). Table~\ref{tab:loops} gives the results under three sets of assumptions: none; \emph{abi}, callees preserve $r_4, \dots, r_7$ and SP, as the calling convention requires; and \emph{abi} with \emph{frame} where the checker allows it (functions that take no address of their frame other than the frame pointer are assumed not to have it written by other code). The checker accepts every certificate, and in the emulated runs no loop takes more back edges per entry than its certified bound. With both assumptions, 376 of the 524 loops of the binaries are certified, and 160 of the 174 functions with loops that receive a certified cycle bound (kernels and division routines) have no trusted row left. In every kernel, in both columns with \emph{abi}, every function keeps the calling convention of Definition~\ref{def:convention}, so that the checker establishes \emph{abi} on the code (Theorem~\ref{thm:machine-conv}): the figures of the column \emph{abi} hold under the sole hypothesis that the saved registers are not overwritten through a pointer (no stack smashing), and those of the last column under that hypothesis and \emph{frame}. Before the producer gave the offsets of the stack pointer for the functions without certified loops, the check failed in 9 kernels for lack of them.

\begin{table}[ht]
\centering
\small
\begin{tabular}{@{}lrrr@{}}
\toprule
 & none & \emph{abi} & \emph{abi}, \emph{frame} \\
Calling convention established on the code (of 29 kernels) & & 29 & 29 \\
\midrule
Loops of the kernels certified (of 403) & 143 & 200 & 260 \\
\quad built with \texttt{-fomit-frame-pointer} & 146 & 204 & 261 \\
\quad certified bound below the source bound & 84 & 116 & 153 \\
\quad no source bound & 59 & 84 & 107 \\
Loops of the division routines certified (of 116) & 116 & 116 & 116 \\
Functions with a certified cycle bound & 116 & 116 & 121 \\
Loop rows still trusted in these bounds & 41 & 35 & 17 \\
\quad built with \texttt{-fomit-frame-pointer} & 39 & 34 & 16 \\
Functions where \emph{frame} is refused (of 265) & & & 26 \\
\bottomrule
\end{tabular}
\caption{Loop bounds certified on the machine code of the 29 TACLeBench kernels (Cortex-M0, Clang \texttt{-O1}). Certificates rejected: none; emulated runs above a certified bound: none. The figures of the other rows and of the text are those of the default build.}
\label{tab:loops}
\end{table}

Of the certified bounds, none is above the source bound (the producer keeps the source row then); those below it are often below by one, since the binary counts back edges and a rotated loop executes its body once before its first back edge. Certified loop bounds also bound loops that the source analysis does not: 9 more functions receive a certified cycle bound than from the source bounds alone, 116 in all (121 with \emph{frame}; among them \texttt{fac\_fac}, \texttt{fir2dim\_pin\_down}, the copy and fill loops of \texttt{md5} and \texttt{pm}, and three functions of \texttt{bitcount}). The checker refuses \emph{frame} in 26 of the 265 functions that the producer analyses (25 without frame pointer), which put another address of their frame in a register, mostly to pass a local array to a callee, where the assumption would indeed not be justified.

\paragraph{Floating point on its machine code.} The Cortex-M0 has no floating-point unit: each operation on a \texttt{float} or a \texttt{double} calls a routine of the run-time library, which the figures above leave without code, so that every function reaching one has no bound. With option \texttt{--rtlib}, each kernel's object is first linked (\texttt{ld -r}) with a real library, which adds the routines the object needs as the toolchain's link would, and these are decoded, bounded, certified and run on their machine code like the kernels' own functions. Two libraries were used: the soft-float routines of LLVM's compiler-rt, at the tag of the Clang used here (18.1.8), compiled by Clang at \texttt{-O2} (no loop, no indirect jump); and the Cortex-M0 build of GCC's libgcc shipped with arm-none-eabi-gcc 13.2.1, unmodified, whose float and double divisions dispatch their special cases through jump tables kept in its read-only data, which the certificate holds with the code (Proposition~\ref{prop:jump-table}). With compiler-rt and both assumptions, 187 of the 267 functions receive a certified cycle bound instead of 121 (181 instead of 116 with no assumption), including the entry functions of 10 kernels instead of 6 (\texttt{complex\_updates}, \texttt{iir}, \texttt{ludcmp} and \texttt{st} added), whose bounds are 1.00 to 5.53 times the observed cycles (5.53 for \texttt{ludcmp} and 3.50 for \texttt{iir}, not analysed further); 256 functions receive a stack bound instead of 151. With libgcc, 187 functions and 10 entry functions are certified (173 and 8 before its tables were followed; \texttt{ludcmp} and \texttt{st} added, with bounds 4.02 and 1.26 times the observed cycles), and 247 functions receive a stack bound. The checker accepts every certificate, and no emulated run exceeds a bound. The calling convention is no longer established everywhere, so \emph{abi} is assumed where the checker says it is not: with compiler-rt in one function (\texttt{ludcmp\_test}, which restores the stack pointer from its frame pointer, now certified because it has a bound), with libgcc in hand-written assembly routines of 13 programs (\texttt{\_\_aeabi\_cfrcmple}, for instance, swaps its operands and branches into another routine). The observed cycles change accordingly: \texttt{st} runs in 6.07 million cycles with compiler-rt and 4.76 million with libgcc, against 1.78 million when the routines are computed outside the processor model. The library's code, like the kernels', is trusted to be the code the program is linked with.

\paragraph{Larger programs.} TACLeBench's 23 sequential and 2 application programs were analysed in the same way (with compiler-rt, \emph{abi} and \emph{frame}), with a limit of one hour each; the debugging output left in two of them at the pinned revision was removed in the work copy. In the 23 programs whose analysis finished (that of \texttt{mpeg2} and \texttt{susan} did not), 316 of 455 functions, the run-time routines included, receive a certified cycle bound, 409 a stack bound, and 614 of 1{,}925 loops are certified on the machine code; the checker accepts every certificate, and no emulated run exceeds a bound. The entry function is certified in two programs, \texttt{lift} and \texttt{statemate}; in the others some function on the call paths keeps a loop that neither the source analysis nor the machine code bounds, calls a routine without code, or is recursive. These programs contain 8 jumps through a table, in both forms of Proposition~\ref{prop:jump-table}; all are recognised, and 3 lie in functions that receive a certificate, among them \texttt{lift}'s, which was the only obstacle to its bound. The calling convention is established on the code of every function of the 23 programs, these 3 included. One program, \texttt{cjpeg\_transupp}, returns a wrong result on the Cortex-M0: it reads past the end of an array, so that its expected result holds only for the layout of the data on the host where it was written.

Table~\ref{tab:steps} records how the counts grew as the checker and the producer were extended, in the order in which the extensions were made. It is not an ablation: no extension was measured without the ones before it, and the counts of a step depend on those before it (the linear programs, for instance, are retried in ways added later). At each step no loop was lost, no certified bound became higher, every certificate was accepted, and no emulated run exceeded a bound. Most of the loops added by the congruences leave on equality with a counter stepped by more than one; the potentials per instruction from a linear program are needed where the counter moves between a register and a word of the frame inside the loop (in \texttt{countnegative} and \texttt{isqrt}); the exception for the frame pointer added 40 loops in 12 kernels (before it, the checker refused \emph{frame} in 190 functions, almost all only because of the prologue \texttt{ADD r7, SP, \#i}); the exact effects added loops of \texttt{bitonic}, \texttt{fft}, \texttt{fir2dim}, \texttt{quicksort} and \texttt{matrix1}, the exact shift loops that depend on a size computed by a shift (in \texttt{fft}, $1 \ll 10$), and it lowered four bounds of \texttt{bitcount} from $2^{31} - 2$ to 255; the disjunctive invariants certified every loop of the division routines (bound 31 back edges, source bound 32) and loops of \texttt{fft}, \texttt{matrix1} and \texttt{pm}; the rounding of sums added six loops of \texttt{bsort}, \texttt{fft} and \texttt{fir2dim} that count a byte offset up to zero; the last step, a change of the trusted link and not of the producer (paragraph ``Another compiler''), matched a binary loop of \texttt{minver} with its source loop, which has no bound, instead of a bounded one, so that the function lost its source bounds and the certified bound of another of its loops, above the source bound before, was kept; following the jump of \texttt{bitcount\_main}, whose index only the invariant bounds, added its three loops, with \emph{frame} only. The constants of the literal pool and the ranges near $2^{32}$ repeatedly made HiGHS fail numerically on linear programs of loops certified before; the producer's retries (Section~\ref{sec:loops}) and its own analysis, which ignores constants above $2^{20}$ and, when a loop of a function is left without a potential, analyses the function again with literal loads and shifts giving any value, were added so that no loop is lost. Tracking the words of the frame had raised the count with \emph{abi} from 94 to 105.

\begin{table}[ht]
\centering
\small
\begin{tabular}{@{}lrrr@{}}
\toprule
Checker and producer & none & \emph{abi} & \emph{abi}, \emph{frame} \\
\midrule
Zones, tracked words of the frame & 73 & 105 & 107 \\
\quad + congruences modulo $2^e$ & 116 & 167 & 169 \\
\quad + potentials per instruction (linear program) & 121 & 173 & 175 \\
\quad + exception for the frame pointer & 121 & 173 & 215 \\
\quad + exact \texttt{NEGS}, \texttt{MVNS}, \texttt{STM}/\texttt{LDM}, literal loads & 131 & 183 & 229 \\
\quad + exact \texttt{LSLS}, exact Farkas multipliers & 136 & 193 & 246 \\
\quad + disjunctive invariants & 136 & 193 & 250 \\
\quad + rounding of sums, wrap thresholds, transported constants & 142 & 199 & 256 \\
\quad + target data model, matching by line ranges & 143 & 200 & 257 \\
\quad + jump indices bounded by invariants, shifts & 143 & 200 & 260 \\
\bottomrule
\end{tabular}
\caption{Loops of the kernels certified (of 403, default build) after each extension, in the order in which they were made; not an ablation. The 116 loops of the division routines are certified from the disjunctive invariants on, under every assumption, and none before.}
\label{tab:steps}
\end{table}

Table~\ref{tab:fail} classifies the 146 loops not certified with both assumptions, all in the kernels, by the first obstacle that the producer meets. Most have a counter that the invariants track but no potential among the producer's candidates; 62 of these leave on equality, and some are in functions that pass an address of their frame to callees or store through pointers, the cause that Sewell et al.\ also name, for the binary alone, for their failures (no aliasing analysis of memory \cite{sewell2016complete}). Others are not linear: \texttt{bitcount} has loops that run once per set bit (\texttt{x \&= x - 1}), and in \texttt{bsort} the number of passes depends on the data.

\begin{table}[ht]
\centering
\small
\begin{tabular}{@{}lr@{}}
\toprule
Cause & Loops \\
\midrule
Counter tracked, no potential among the producer's candidates & 86 \\
\quad of which the loop leaves on equality & 62 \\
Compared value not tracked & 21 \\
Offset of SP unknown & 20 \\
Flags set by an instruction whose effect is not modelled & 12 \\
Several of these & 7 \\
\midrule
In all & 146 \\
\bottomrule
\end{tabular}
\caption{Loops of the kernels not certified with both assumptions, by cause (the first obstacle met by the producer's analysis; a diagnosis, not a proof that no certificate exists).}
\label{tab:fail}
\end{table}

The specification of effects, flags, conditions and stack effects is compared with our interpreter on every 16-bit encoding, in three random states each, with and without the assumption on callees: 297{,}244 executions, without difference (\lean{benchmarks/cycles/test\_effect.py}). The checker refuses certificates altered in their bounds, scales, potentials, invariants, kinds, multipliers, tracked words, offsets of SP, congruences, assumptions or code, and certificates of small loops for each extension of Table~\ref{tab:steps} when the instruction it models is changed: a literal word, the operand of \texttt{NEGS} or \texttt{MVNS}, an \texttt{STM} replaced by an \texttt{LDM}, the amount or operand of an \texttt{LSLS}, an overflowing shift, \texttt{LSRS} for \texttt{LSLS}, a shift by zero that breaks a congruence, the step or start of a counter that leaves on its wrap-around, or its alternative disjuncts and targets removed or changed (\lean{test\_binloops.py}). A run of all kernels with loop certificates takes 7 to 19~minutes with four processes; before the producer closed its zones incrementally (a change that leaves its certificates identical), the slowest took 103~minutes. Choosing, among the integer and the exact fractional dual multipliers of the IPET program, those that give the smallest bound (rows of certified loops with $D > 1$ have fractional multipliers) lowered 10 certified cycle bounds, by up to 8\%, when it was introduced (with the congruences).

\section{Mechanization}
\label{sec:mechanization}

The modules of this paper (\lean{CyclesCore}, \lean{Cycles}, \lean{ThumbCore}, \lean{Thumb}, \lean{ThumbDataCore}, \lean{ThumbData}, \lean{ThumbConv}) have about 6{,}500 lines of Lean~4 and 257 theorems, in the repository of Paper~I; the executable parts (\lean{\dots Core}) do not use mathlib and are compiled into the checker program. Table~\ref{tab:lean3} lists which Lean statements back the paper's. Each module's examples check a certificate of a counting loop by \lean{decide} and refuse a bound one lower.

\begin{table}[ht]
\centering
\small
\begin{tabular}{@{}>{\raggedright\arraybackslash}p{0.27\textwidth}>{\raggedright\arraybackslash}p{0.5\textwidth}>{\raggedright\arraybackslash}p{0.15\textwidth}@{}}
\toprule
Paper statement & Lean & Module \\
\midrule
Def.~\ref{def:cycle-certificate}, Prop.~\ref{prop:cycles} & \lean{checkFn}, \lean{checkProgram}, \lean{Run}, \lean{cycles\_\allowbreak sound}, \lean{program\_\allowbreak sound} & \lean{Cycles} \\
Def.~\ref{def:machine-run}, Thm.~\ref{thm:machine-cycles} & \lean{decode}, \lean{trans}, \lean{buildFn}, \lean{checkBinary}, \lean{MRun}, \lean{translate}, \lean{binary\_\allowbreak sound} & \lean{Thumb} \\
Def.~\ref{def:data-run} & \lean{effect}, \lean{stackEffect}, \lean{flagsAdd}, \lean{flagsSub}, \lean{condHolds}, \lean{EffOk}, \lean{StackOk}, \lean{DOk}, \lean{DRun} & \lean{ThumbData} \\
Def.~\ref{def:loop-certificate} & \lean{DFun}, \lean{LoopCert}, \lean{loopRow}, \lean{obligations}, \lean{congOk}, \lean{checkData}; disjunctive invariants \lean{invsAt}, \lean{tgtIdx}, \lean{tgtInv}; example \lean{exWrap\_\allowbreak sound} (a counter that leaves on its wrap-around), \lean{exWrap\_\allowbreak needs\_\allowbreak alt} & \lean{ThumbData} \\
Lem.~\ref{lem:cond-facts} & \lean{condFacts\_\allowbreak sound}, \lean{sub\_\allowbreak nv}, \lean{sub\_\allowbreak z}, \lean{add\_\allowbreak z} & \lean{ThumbData} \\
Congruences, rounding & \lean{congPE\_\allowbreak sound}, \lean{gExp\_\allowbreak dvd}, \lean{tightenL\_\allowbreak holds}, \lean{gExpS\_\allowbreak dvd}, \lean{tightenS\_\allowbreak holds}, \lean{condFactsM\_\allowbreak sound}; example \lean{exZero\_\allowbreak sound} (a counter that leaves when it wraps around to zero), \lean{exZero\_\allowbreak needs\_\allowbreak cong}; example \lean{exStep\_\allowbreak sound} (a counter stepped by 4, exit on equality), \lean{exStep\_\allowbreak needs\_\allowbreak cong} & \lean{ThumbData} \\
Thm.~\ref{thm:machine-loops} & \lean{postS\_\allowbreak eval}, \lean{step\_\allowbreak ok}, \lean{chain\_\allowbreak ok}, \lean{chain\_\allowbreak rows}, \lean{drun\_\allowbreak mrun}, \lean{data\_\allowbreak sound}; examples \lean{exData\_\allowbreak sound}, \lean{exData\_\allowbreak run} (a run of 68 cycles: the bound is attained) & \lean{ThumbData} \\
Prop.~\ref{prop:jump-table} & \lean{JT}, \lean{transJ}, \lean{xop}, \lean{xres}, \lean{pcWrite}, \lean{pcDestV}, \lean{DOkM}, \lean{XOk}, \lean{wstep}, \lean{wrun}, \lean{jumpOk}, \lean{interiorOk}, \lean{wpath}, \lean{wpred}, \lean{wrun\_\allowbreak sound}, \lean{back\_\allowbreak walk}, \lean{jump\_\allowbreak good}, \lean{steps\_\allowbreak good}; second form: \lean{rngOps}, \lean{jumpRng}, \lean{interiorR}, \lean{back\_\allowbreak walkR}, \lean{jump\_\allowbreak good\_\allowbreak rng}; third form: \lean{invBound}, \lean{congFits}, \lean{jumpInv}, \lean{inv\_\allowbreak prefix}, \lean{jump\_\allowbreak good\_\allowbreak inv}, the congruence of a shift \lean{shlCong\_\allowbreak sound}; examples \lean{exSw\_\allowbreak sound} (a \texttt{switch} of three cases), \lean{exLs\_\allowbreak sound} (with \texttt{BLS}), \lean{exSw\_\allowbreak target\_\allowbreak missing}, \lean{exMask\_\allowbreak sound} (no \texttt{CMP}: \texttt{switch (x \& 3)}), \lean{exMask\_\allowbreak wide}, \lean{exPath\_\allowbreak sound} (an index set on two paths, bounded by the invariant), \lean{exPath\_\allowbreak wide} & \lean{ThumbData} \\
Defs.~\ref{def:linked-run}, \ref{def:convention}, Thm.~\ref{thm:machine-conv} & \lean{DRunL}, \lean{PushPopOk}, \lean{SaveSafe}, \lean{convOk}, \lean{convAll}, \lean{ConvInv}, \lean{conv\_\allowbreak step}, \lean{conv\_\allowbreak ret}, \lean{drunL\_\allowbreak drun}, \lean{data\_\allowbreak sound\_\allowbreak conv}; examples \lean{exConv\_\allowbreak run} (a linked run of 68 cycles), \lean{exSave\_\allowbreak conv}, \lean{exClobber\_\allowbreak refused}, \lean{exSwConv\_\allowbreak sound} (a switch), \lean{exSwClobber\_\allowbreak refused} (a register written only through the table) & \lean{ThumbConv} \\
\bottomrule
\end{tabular}
\caption{Paper statements and their Lean counterparts (namespace \lean{ImperativeOptimization.Thumb}, \lean{Cycles} for the first row).}
\label{tab:lean3}
\end{table}

Not mechanized: that the specification (decoder, cycles, effects on registers, flags, stack pointer and memory, conditions, the destinations of \texttt{ADD PC} and \texttt{MOV PC}, the exact values of \texttt{ADD rd, PC}, \texttt{ADR} and the loads from the code; zero-wait-state memory, no interrupts) describes the Cortex-M0, that the code in a certificate is that of the binary and is not written during a run, that \texttt{PUSH} and \texttt{POP} store and load as stated in Definition~\ref{def:linked-run}, that callees preserve $r_4, \dots, r_7$ and SP when this is assumed (\emph{abi}) without being checked on their code, that no store through a pointer and no callee writes the words where a function saves registers when the convention is checked (hypothesis~(3) of Definition~\ref{def:linked-run}), that no other code writes a function's frame when this is assumed (\emph{frame}), and that the trusted rows hold. The hypotheses of Theorem~\ref{thm:machine-loops} are satisfiable: a run of the example's code in the data semantics, of exactly 68 cycles, is built in Lean by a concrete simulator whose acceptance implies the relations of Definition~\ref{def:data-run} (\lean{exData\_\allowbreak run}), so the certified bound is also attained; the same run is a linked data run (\lean{exConv\_\allowbreak run}), so the hypotheses of Theorem~\ref{thm:machine-conv} are satisfiable too.

\section{Limitations}
\label{sec:limitations}

The specification of the processor is written for the purpose and compared with an interpreter written by the same hand, not with hardware; a model of the instruction set validated on silicon \cite{fox2010trustworthy} would make it more trustworthy. Memory is modelled only for the words of the frame accessed through SP; loops whose counters are kept elsewhere in memory, or in a frame whose address is taken other than by the frame pointer, are not certified on the machine code; neither are loops whose tests depend on instructions whose effects are approximated. Potentials are differences of at most two registers with per-instruction constants. Floating point is covered only with a run-time library linked. Jumps through a table are followed only in the three forms of Proposition~\ref{prop:jump-table}; other indirect jumps, calls through pointers and tables whose index the producer's zones and congruences do not bound within 1024 values are refused. The bounds are for one processor, the Cortex-M0 with zero-wait-state memory, and without interrupts. The kernels are small programs.

\section*{Conclusions}

Cycle bounds of machine code can be certified by a checker proved sound in Lean, with the graph built by the checker from the code, so that the trusted base is a specification of the processor. The loop bounds of the IPET program can be certified on the same code, for the loops whose counters live in registers or in the words of the frame, which removes the compiler, the debug information and the source analysis from the trusted base for those loops. The calling convention of the callees need not be assumed either: the checker establishes it on their code, under the hypothesis that no store through a pointer overwrites the saved registers. Linked with a run-time library, floating-point code is bounded on the machine code of its routines. Jumps through a table are followed when a range check, an instruction of the window or the certified invariant bounds their index, their targets established by evaluating the window on the code and proved from the data semantics. Future work is a treatment of addresses of the frame passed to callees, which would let \emph{frame} be checked in the same way, more exact effects, and a specification of the processor validated on hardware.

\section*{Artifact}

The Lean modules of Table~\ref{tab:lean3}, the analyses of the binaries, the interpreter, the producers of certificates and the scripts that produce the measurements of this paper are in release v1.0 of the repository \url{https://github.com/ldallolio/crosscheck}, archived at Zenodo, \href{https://doi.org/10.5281/zenodo.24425070}{doi:10.5281/zenodo.24425070}, and at Software Heritage, \href{https://archive.softwareheritage.org/swh:1:rel:70c77563037e6db2a55b48b3ad4963c86a67773f}{\nolinkurl{swh:1:rel:70c77563037e6db2a55b48b3ad4963c86a67773f}}.

\section*{Statement on the Use of AI}

Following the recommendations on the responsible release of AI-generated mathematics \cite{agmai2026responsible}, we document how AI was used in this work. In November 2026 an AI coding assistant (Anthropic's Claude, model Opus~5.5, and Haiku~5.5 inside the component that reads web pages, used through the Claude Code tool) wrote the Lean modules of Table~\ref{tab:lean3}, the cycle and stack analyses of the binaries, the interpreter, the producers of certificates, the OTAWA experiments, the check of time budgets, the tests, ran the measurements and drafted the text of this paper, which takes over sections of the two companion papers; the prompts and a summary of the changes are recorded in the file \lean{AI\_USAGE.md} of the repository. Every Lean proof is checked by the Lean kernel. I have read and verified the text and take full responsibility for the content. The AI assistant checked the attributions to \cite{cullmann2007data,crary2000resource,sewell2016complete,granger1989static} against the full texts of these papers, which the author supplied.

\bibliographystyle{unsrt}
\bibliography{../crosscheck}

\end{document}
